% Run GCT_KRON_20261120; author/checker GPT-6 Astra Pro, effort not exposed.
% Source scope updated after official cote133 facsimile verification; see A-029.
\section{Audit of the archive: conventions and the positivity gate}
The supplied sources are transcriptions, including explicit uncertain and
deleted readings. Every attached archive page has now been compared visually
with its official Montpellier facsimile: cote133 scans006,008,009,019,028;
cote2 scan017; cote12 scan041; cote43 scan037; cote119 scan217; and the
scope-exclusion pages cote4 scan007, cote16 scan038, cote79 scan051,
cote89 scan018. Nine additional cote133 pages020--027,029 were also read.
The source manifests and check events record the exact scope. An uncertain
surname on cote133 scan008 remains uncertified and is not used in any proof.
The statements below retain their explicit mathematical interpretations;
a deleted sentence or a question is not turned into an endorsed theorem.
The secondary study SA17 remains distinguished from the original pages,
and every relevant assertion, definition, question and scope exclusion is
recorded in the inventory.

For a finite group $G$, $B(G)$ is the group completion of finite $G$-sets under
disjoint union, with product induced by Cartesian product.  We distinguish
\[
 p_r(E)=[\{A\subset E:|A|=r\}],\qquad
 \sigma_r(E)=[\operatorname{Sym}^rE],\qquad
 \lambda_t^B(E)=\sigma_{-t}(E)^{-1}.
\]
Both $p_t$ and $\lambda_t^B$ are pre-$\lambda$ structures, but they are different.
The second is the conventional Burnside structure compatible with exterior
powers after linearization.  A ``special $\lambda$-ring'' means that the
universal product and iteration identities also hold (SGA 6, V, Definition
2.4, p.~320).  This is stronger than the addition identity for a pre-structure.

\subsection*{A-001. The subset reading of scan 019 contradicts its boxed identity}
\textbf{Source and status.} Cote 133, scan 019, lines 4--10 and 43--54;
SA17, lines 491--496.  Under the study's interpretation of $\mathcal P_r(E)$
as $r$-element subsets, the boxed identities are false in $B(G)$ and after
linearization.  The corrected identities are
\[
 \sigma_2(E)=p_2(E)+E,\qquad
 \lambda_2^B(E)=E^2-\sigma_2(E).
\]
\textbf{Proof.} A multiset of size two has either one repeated member or two
distinct members.  This gives the first equivariant bijection.  Taking the
coefficient of $t^2$ in $\sigma_t\lambda^B_{-t}=1$ gives the second equality.
Let $G=C_2$ and let $E$ be its free orbit.  Then $E^2=2E$ and $p_2(E)=1$.
Thus the proposed equality $2p_2(E)+E=E^2$ is $2+E=2E$.  Applying the number
of $G$-fixed points gives $2=0$, a contradiction.  Under linearization over
$\mathbb Q$, $E$ becomes $1+\varepsilon$, where $\varepsilon$ is the sign
character.  The two sides become $3+\varepsilon$ and $2+2\varepsilon$,
and are distinct.  In contrast, $\lambda_2^B(E)=E-1$ linearizes to
$\varepsilon=\bigwedge^2(1+\varepsilon)$.
The cardinality equality $2\binom d2+d=d^2$ is correct; forgetting the group
action conceals the error.  This also refutes SA17's assertion that the
Burnside identity was verified in degree two. \hfill$\square$

\paragraph{Dictionary.}
Finite $G$-sets give finite etale schemes with action of the constant group
scheme $G_{\mathbb Z}$ and permutation locally free sheaves on
$BG_{\mathbb Z}$.  The subset construction forgets the signs on an exterior
basis: a stabilizer can permute a chosen subset with odd sign.  The exterior
power is therefore not the permutation sheaf on subsets.  Constant group
schemes are constructed in SGA 3, I, 4.1, original p.~23; the representation
and comodule dictionary is I, Proposition 4.7.2, pp.~34--35.

\subsection*{A-002. Canonical Burnside operations are not special}
\textbf{Source and status.} The question $\lambda^i(xy)=?$ in cote 133,
scan 019, line 50 has a negative answer for the canonical inverse-symmetric
operations.  The general non-specialness and the signed subset formula are
known results; see Bouc--R\"okaeus, \emph{A note on the $\lambda$-structure on
the Burnside ring}, Theorem 1.1, formula (1.2).  We give a short explicit
counterexample to locate the failure relevant to the archive.

\textbf{Proof of the counterexample.}
Take $G=C_2\times C_2$, let $H_1,H_2,H_3$ be its three subgroups of order
two, put $X_i=G/H_i$, and put $E=G/1$.  Orbit counting gives
\[
 X_i^2=2X_i,\qquad X_iX_j=E\quad(i\ne j),\qquad
 \lambda_2^B(X_i)=X_i-1.
\]
If $H\leq G$, a multiset is $H$-fixed exactly when its multiplicity is
constant on each $H$-orbit.  For the mark $m_H$ this proves
\[
 m_H(\sigma_t(S))=\prod_{O\in H\backslash S}(1-t^{|O|})^{-1},\qquad
 m_H(\lambda_t^B(S))=\prod_{O\in H\backslash S}(1-(-t)^{|O|}).
\]
Consequently $\lambda_2^B(E)$ has marks $0,-2,-2,-2,6$ at
$G,H_1,H_2,H_3,1$, respectively, and
\[
 \lambda_2^B(E)=3E-X_1-X_2-X_3.
\]
The special degree-two identity would instead give
\[
 \lambda^2(X_1X_2)
 =X_1^2\lambda^2(X_2)+X_2^2\lambda^2(X_1)
       -2\lambda^2(X_1)\lambda^2(X_2)=2E-2.
\]
Their difference $E-X_1-X_2-X_3+2$ has $G$-mark $2$, hence is nonzero.
This difference linearizes to zero, illustrating why a character calculation
alone cannot verify specialness of the Burnside structure. \hfill$\square$

For completeness the subset pre-structure also fails the special
product identity. For the free two-point $C_2$-set $E$ one has
$p_2(E)=1$ and $E^2=2E$. The disjoint-union subset formula gives
$p_2(E^2)=p_2(2E)=2+E^2=2+2E$, whereas the special product
polynomial gives $E^2+E^2-2=4E-2$. Their difference has
$C_2$-fixed-point mark four. Both subset and canonical operations
do satisfy the pre-structure addition identity: subsets of a
disjoint union split by their two sizes, and the same argument
holds for multisets before inverting their total series.

\paragraph{The correct signed formula.}
Bouc--R\"okaeus, Theorem 1.1 (1.2), states
\[
 \lambda_n^B(S)=(-1)^n\sum_{\mu\vdash n}(-1)^{\ell(\mu)}
 \frac{\ell(\mu)!}{\prod_j\alpha_j(\mu)!}[\mathcal P_\mu(S)],
\]
where $\mathcal P_\mu(S)$ is the $G$-set of ordered pairwise disjoint subsets
with sizes given by the parts of $\mu$, and $\alpha_j(\mu)$ counts repeated
parts.  It is a formula in the group completion, not a positive formula.
Linearization commutes with these operations because multisets form a basis
of symmetric powers of a permutation module and the symmetric/exterior
generating-series identity holds in the representation ring.  The special
structure on the latter is SGA 6, VI, Theorem 3.3, pp.~392--393; it does not
transfer specialness backwards across a noninjective homomorphism.

\subsection*{A-003. Effective permutation lifts cannot supply nontrivial irreducibles}
\textbf{Claim.} For a field of characteristic zero and a finite group $G$,
no nontrivial irreducible representation is the linearization of an effective
finite $G$-set.  In particular, a nontrivial Specht representation has no
effective finite symmetric-group-set lift.

\textbf{Proof.} For a finite $G$-set $X$, the invariant subspace of $k[X]$
has basis the sums of the elements of each orbit.  Hence its dimension is
the number of orbits.  If $X$ is nonempty this dimension is positive.  A
nontrivial irreducible representation has no nonzero invariants; if it did,
its invariant subspace would be the whole representation and the
representation would be trivial.  The empty set linearizes to zero.
These exhaust the effective possibilities. \hfill$\square$

\paragraph{Consequence and dictionary.}
Virtual permutation resolutions, including Jacobi--Trudi constructions from
Young subgroups, can still yield exact signed Kronecker formulas.  This
argument rules out deriving a positive irreducible-basis formula merely by
treating each Specht class as an effective permutation class.  It does not
rule out an independent positive rule or a geometric realization by
cohomology.  On $BG$, the invariant global sections of a nonempty
permutation bundle are nonzero, whereas a nontrivial irreducible bundle has
none.  The group-completion step is the precise loss of effectiveness.

\subsection*{A-004. The monomial Cauchy coefficients on scan 028 need not be effective}
\textbf{Source and status.} Cote 133, scan 028, lines 16--24 asks whether
the summands in the displayed expansion are actual direct summands when
$\xi$ is effective.  Under its explicitly given monomial symmetric
coefficient convention, the answer is no, already over $\mathbb Q$.

\textbf{Proof.} In degree two the displayed expansion is
\[
 \tau^2(E)=\lambda^2(E)\rho_1^2+
             (E^2-2\lambda^2(E))\rho_2.
\]
Here $\rho_1^2$ is the regular representation of $S_2$ and $\rho_2$ is its
trivial representation; $E^2-2\lambda^2(E)=\psi^2(E)$.
Let $E$ be the two-dimensional standard representation of $S_3$, regarded
as a bundle on $BS_3$.  Its characters on the classes of the identity,
transpositions and three-cycles are $(2,0,-1)$.  Thus
\[
 \bigwedge^2E=\varepsilon,\quad
 \operatorname{Sym}^2E=1+E,\quad
 \psi^2(E)=1+E-\varepsilon.
\]
The coefficient of $\varepsilon$ in the last expression is negative, so it
is not the class of an effective representation.  In contrast the actual
$S_3\times S_2$ decomposition is
\[
 E^{\otimes2}=(1+E)\boxtimes1+\varepsilon\boxtimes\operatorname{sgn}.
\]
Expanding $\rho_1^2=1+\operatorname{sgn}$ cancels the negative coefficient
and recovers this genuine decomposition.  The monomial expansion is a
valid virtual identity in this example, but not a decomposition into
effective terms. \hfill$\square$

\paragraph{Dictionary.}
The example is a locally free sheaf on the ringed topos $BS_{3,\mathbb Q}$,
so it directly lies in the scope of scan 028.  Schur functors provide the
characteristic-zero direct-summand decomposition of a tensor power; changing
from Schur coefficients to monomial symmetric coefficients destroys
termwise effectiveness.  A monomial Cauchy identity therefore gives no
saturation statement.

\subsection*{A-005. The split quadratic Grothendieck group on scan 006 need not be special}
\textbf{Source and distinction.} Cote 133, scan 006, lines 6--10 defines
the group completion of the isometry monoid, without imposing relations
from metabolic spaces.  Lines 20--32 ask whether its exterior operations
are special when $2$ is invertible.  For this literal split definition the
answer is no.  This does not contradict the theorem for the modern
Grothendieck--Witt group: Zibrowius, \emph{Symmetric representation rings
are $\lambda$-rings}, introduction theorem, p.~1056, uses Definition 1.2,
p.~1060, which also imposes metabolic relations.

\textbf{Proof.} Work on the ringed classifying topos $BC_3$ over
$k=\mathbb F_3$.  Write $1,a,b$ for the classes of the indecomposable
modules $J_1,J_2,J_3$ on which a generator acts by one unipotent Jordan
block.  Since $k[C_3]=k[t]/(t^3)$, Jordan normal form implies that its
split representation ring is free as an abelian group on $1,a,b$.
The tensor and exterior identities in that ring are
\[
 a^2=b+1,\qquad ab=2b,\qquad b^2=3b,
 \qquad \lambda^2(a)=1,\quad\lambda^2(b)=b,\quad\lambda^3(b)=1.
\]
For completeness, on $J_2\otimes J_2$ the nilpotent operator $g-1$ has
ranks $2,1,0$ in powers one, two, three, so its blocks are $J_3+J_1$.
The module $J_3=k[C_3]$ is regular.  The map
$g\otimes v\mapsto g\otimes g^{-1}v$ identifies its tensor product with
$V$ under diagonal action with $\dim(V)$ copies of the regular module.
This proves the other tensor identities.  The determinants of unipotent
blocks are one; $J_3$ is self-dual and
$\bigwedge^2J_3\simeq J_3^*\otimes\det J_3\simeq J_3$, proving the
exterior identities.

In any pre-$\lambda$-ring the third Newton operation is
\[
 \psi^3(z)=z^3-3z\lambda^2(z)+3\lambda^3(z).
\]
It is additive, as follows by comparing coefficients of the formal
logarithmic derivative of $\lambda_t(z+w)=\lambda_t(z)\lambda_t(w)$.
The preceding identities give
\[
 \psi^3(a)=2b-2a,\qquad \psi^3(b)=3.
\]
Now take the symmetric hyperbolic modules $x=H(J_2)$ and $y=H(J_3)$.
Since $2$ is invertible, symmetric bilinear forms and quadratic forms
are equivalent.  Forgetting a form defines a ring homomorphism that
commutes with exterior powers from the split quadratic group to the
split representation ring.  The images of $x,y,xy$ are $2a,2b,8b$.
Consequently the images of $\psi^3(xy)$ and $\psi^3(x)\psi^3(y)$ are
respectively
\[
 24\cdot1\quad\hbox{and}\quad24b-24a.
\]
They are unequal in the free abelian group on $1,a,b$.  In a special
$\lambda$-ring every Newton operation is multiplicative: the universal
product identities imply this polynomial identity, which can be checked
on formal sums of line elements, where $\psi^3(\sum_i l_i)=\sum_i l_i^3$.
Thus the split quadratic group in this example is not special.
\hfill$\square$
\paragraph{Dictionary and significance.}
This is a counterexample on a classifying ringed topos, exactly the
framework of scan 006.  Quotienting the split group by metabolic relations
changes the mathematical object.  It is therefore essential to specify
the exact category and relations when importing a modern positivity or
$\lambda$-ring theorem into the manuscript's notation.  The counterexample
does not address a semisimple characteristic-zero restriction of the
split group; it refutes the stated generality with $2$ invertible.

\subsection*{A-006. Without invertibility of two, strict nonsingularity is not preserved by exterior powers}
\textbf{Source and status.} Cote 133, scan 006, lines 12--19 claims a
pre-$\lambda$ structure before assuming that $2$ is invertible and states
that ranks are even when $2$ is not invertible.  The rank assertion is
correct for a connected ringed topos with locally free modules of
constant rank; the claimed exterior-power pre-structure is not.

\textbf{Proof.} The polar form of a quadratic form has diagonal entries
$2q(e_i)$.  Modulo two its matrix is alternating.  An alternating matrix
of odd size has determinant zero over every ring, by the integral
alternating determinant identity.  Therefore the determinant of an odd
rank polar matrix belongs to the ideal $(2)$.  If it is a unit, then
$2$ is a unit.  Applied locally and then glued, this proves the even-rank
assertion whenever $2$ is not globally invertible and the rank is
constant.

Over a field of characteristic two the hyperbolic plane is a strictly
nonsingular quadratic module of rank two.  Its second exterior power
has rank one.  A rank-one quadratic module in characteristic two has
zero polar form, so cannot be strictly nonsingular.  Thus exterior
powers cannot define the claimed pre-structure on this collection of
objects.  When $2$ is invertible, exterior powers of the associated
nonsingular symmetric forms do give a pre-$\lambda$ structure, and
their ranks are binomial coefficients.  Specialness remains false for
the split group by A-005. \hfill$\square$

\subsection*{A-007. Exterior operations do not descend to the Witt quotient}
\textbf{Source and status.} Cote 133, scan 008, lines 32--38 asks whether
the pre-$\lambda$ structure descends modulo hyperbolic modules.  It does
not, even over a field of characteristic different from two.

\textbf{Proof.} The hyperbolic plane $H=\langle1,-1\rangle$ is zero in
the Witt group, while
\[
 \lambda^2(H)=\langle-1\rangle
\]
has odd rank and is nonzero by the dimension-modulo-two homomorphism.
Thus the hyperbolic ideal is not stable under exterior operations.
For $2$ invertible it is indeed an ideal: for a nonsingular symmetric
module $(M,b)$, the identification $M\simeq M^*$ induced by $b$ gives
\[
 H(E)\otimes(M,b)\simeq H(E\otimes M).
\]
Taking products with group-completed classes preserves the image of
the hyperbolic homomorphism.  Rank modulo two descends because
$\operatorname{rank}H(E)=2\operatorname{rank}E$; integral rank does
not descend, since $H(k)$ has rank two but is zero in the quotient.
\hfill$\square$

\subsection*{A-008. The general topos caveat on scan 008 is essential}
\textbf{Source and status.} Cote 133, scan 008, lines 2--8 asserts that
$q\perp(-q)$ is hyperbolic for modules over a ring, with a caveat for
general ringed topoi.  Lines 11--20 suggest a stable-hyperbolic
description of equality in the quotient.  The ring assertion holds for
finite projective modules with nonsingular polar form; the unrestricted
topos description fails in characteristic two, including its proposed
stable criterion.

\textbf{Proof over a ring.} Let $M$ be finite projective and let $b$ be
the nonsingular polar form of $q$.  The surjection
$M^*\otimes M^*\longrightarrow\operatorname{Sym}^2M^*$ splits as a
map of modules because its target is finite projective.  Hence choose
a bilinear form $c$ with $c(u,u)=q(u)$, and use $b$ to choose an
endomorphism $a$ satisfying $b(u,a(u))=q(u)$.  In the coordinates
$u=x-y,v=y$ one has
\[
 q(x)-q(y)=q(u)+b(u,v)=b(u,v+a(u)).
\]
The change of coordinates is invertible, and the last form is the
hyperbolic form after $b$ identifies the second copy of $M$ with $M^*$.
If $2$ is invertible, the choice $c=b/2$ is canonical and the argument
works on every ringed topos.  Without that hypothesis the global choice
of $c$ can fail. 

\textbf{Counterexample on a classifying topos.}
Let $k$ be algebraically closed of characteristic two, and let
$G=\mathbb G_m\rtimes C_2$, where the nontrivial element $s$ acts on
$\mathbb G_m$ by inversion.  Let $V=ke\oplus kf$ have torus weights
$1,-1$, with $s$ exchanging $e,f$, and set $q(xe+yf)=xy$.
This is a nonsingular invariant quadratic form, denoted $Q$.
In fact $G$ is its split orthogonal group: a matrix preserving $xy$
has $ac=bd=0$ and $ad+bc=1$, and hence lies in the diagonal or
anti-diagonal component, each isomorphic to $\mathbb G_m$.
Every $G$-endomorphism of $V$ is scalar, and $V$ is simple: a torus
stable line has one of the two weights and is not stable under $s$.
Thus a $G$-linear copy of $V$ in $V\oplus V$ has the form
$v\mapsto(av,bv)$, and the restriction of $q\perp q$ is
$(a+b)^2q$.  The only isotropic copy is the diagonal.  In contrast a
hyperbolic module has two complementary invariant isotropic submodules.
If $H(E)$ had underlying representation $V\oplus V$, semisimplicity
of this particular representation forces $E\simeq V$.  Hence
$Q\perp(-Q)=Q\perp Q$ is not hyperbolic.

We now prove failure even of the displayed stable criterion.  In any
$G$-representation the weight pair $\{1,-1\}$ is a canonical direct
summand, and an invariant nonsingular form restricts to a nonsingular
form there.  Write this summand as $V\otimes k^r$.  Torus invariance
forces its quadratic form to be $x^TAy$, where $A$ is an invertible
$r\times r$ matrix; $s$-invariance forces $A$ to be symmetric.  Changes
of equivariant coordinates act on $A$ by congruence.  The restriction
of a hyperbolic module $H(E)$ to this weight pair gives an alternating
matrix: if the multiplicity of $V$ in this pair of $E$ is $r$, this
matrix is congruent to
$\left(\begin{smallmatrix}0&I_r\\I_r&0\end{smallmatrix}\right)$.
The matrix for $2Q$ is $I_2$, which is nonalternating.  Adding an
alternating matrix preserves nonalternation.  Consequently no
\[
 2Q\perp H(F)\simeq H(E)
\]
can exist, for any representations $E,F$.

Nevertheless $3Q\simeq Q\perp H(V)$.  Indeed, for the bilinear form
with matrix $I_3$, the basis
\[
 (1,1,1),\quad(1,1,0),\quad(0,1,1)
\]
has Gram matrix $1\oplus
\left(\begin{smallmatrix}0&1\\1&0\end{smallmatrix}\right)$ in
characteristic two.  Thus in the group completion
$2[Q]=[H(V)]$, so $2Q$ is zero in the quotient by the hyperbolic image,
although it is not stably hyperbolic in the displayed sense.
This disproves the unrestricted criterion on scan 008.  Over a field
with no group action, the ordinary Witt cancellation statement quoted
there is compatible with the ring argument; the counterexample concerns
equivariant objects on a classifying topos. \hfill$\square$

\subsection*{A-009. Correct interpretation of the Stiefel--Whitney statements}
\textbf{Source.} Cote 133, scan 009, lines 2--39; SA17, lines 574--575.
The existence, naturality, normalization on lines, and orthogonal-sum
formula are known statements in the etale setting for schemes over
$\mathbb Z[1/2]$. Exact primary source: Esnault--Kahn--Viehweg,
\emph{Coverings with odd ramification and Stiefel--Whitney classes},
\S5.1--5.3, pp.~179--182. The hyperbolic formula is the Proposition in
\S5.5, pp.~183--184. Thus the formula called ``new in this reading''
by SA17 was already in that 1993 source.

The target of the total class on group-completed objects is
$1+\prod_{i>0}H^i_{\mathrm{et}}(X,\mathbb F_2)$, with multiplication of
series. A virtual object can have infinitely many nonzero classes,
because total classes of negatives are formal inverses. The degree-zero
component is one. For ordinary bundles the classes above the rank vanish.
The crossed-out assertion $w\circ H=1$ is false. The cited formula is
\[
 w(H(E))=\sum_{i=0}^{r}(1+\alpha)^{r-i}\bar c_i(E),
 \qquad \alpha=(-1)\in H^1_{\mathrm{et}}(X,\mathbb F_2),
\]
where $\bar c_i$ is the Chern class reduced modulo two in degree $2i$.
For a line bundle, $w_2(H(L))=\bar c_1(L)$. For instance on
$\mathbb P^1_{\mathbb C}$ with $L=\mathcal O(1)$ this is nonzero.

\textbf{Verification of the universal formulas used in the archive.}
Use the cohomologically injective orthogonal splitting principle
(Zibrowius, \emph{The gamma-filtration on the Witt ring of a scheme},
Proposition 3.3, pp.~19--20 of arXiv:1411.3198). If the first
Stiefel--Whitney roots of $M,N$ are $u_i,v_j$, respectively, line-bundle
normalization and the tensor product of the corresponding double-cover
torsors give
\[
 w(M\otimes N)=\prod_{i,j}(1+u_i+v_j),\qquad
 w(\bigwedge^sM)=\prod_{|I|=s}\left(1+\sum_{i\in I}u_i\right).
\]
These are symmetric polynomials in each root alphabet. The fundamental
theorem on symmetric polynomials expresses them universally in the
elementary symmetric polynomials, namely the $w_i$, and the numbers of
roots, namely the ranks. Injectivity of the splitting pullback proves
the original identities. The target of the asserted $\lambda$
homomorphism must be endowed with these Chern-root operations, not with
componentwise cup product as its ring product. The manuscript's
shorthand about an arbitrary strictly local ringed topos does not state
enough topology or splitting hypotheses to identify an equally broad
theorem; the verified statement here is the etale scheme formulation.

\subsection*{A-010. An integral proof of the monomial Cauchy identity}
\textbf{Source and conventions.} Cote 133, scan 028, lines 2--18 and
26--48. Let $K_0(X)$ mean the exact Grothendieck group of locally free
modules, with relations from short exact sequences; use bounded finite
ranks initially; A-017 below removes the boundedness restriction.
On $\bigoplus_{n\geq0}K_0(B_{S_nX})$ use induction along
$S_i\times S_j\subset S_{i+j}$ for the graded product. With these
conventions the displayed identity is integral; it does not require
division by $n!$. On an arbitrary ringed topos, ``bounded finite ranks''
includes all constant-rank modules and their finite virtual combinations.
This is enough for every representation considered in this work.

\textbf{Proof.} Tensor powers of a short exact sequence of locally
free modules have a filtration, preserved by $S_n$, whose associated
graded is the tensor power of the direct sum of the submodule and
quotient. Filter by the number of tensor factors lying in the submodule.
Its quotients are direct sums indexed by subsets of the tensor positions
and are locally free. Consequently
\[
 \tau_t([E]+[F])=\tau_t([E])\tau_t([F]),
 \quad \tau_n([E])=[E^{\otimes n}],
\]
where the right-hand product uses induction. This defines $\tau_t$
on $K_0(X)$ by group completion, with values in units of the completed
graded ring.

For an actual module of constant rank $d$, check the identity
universally on the locally ringed classifying topos of
$GL_{d,\mathbb Z}\times S_n$. The general bundle and its tensor
permutation action are obtained from this universal bundle by its frame
torsor and base change. Pull back to the full flag bundle of the
$GL_d$ bundle. Each successive projective-bundle pullback on $K_0$
is injective by SGA 6, VI, Theorem 1.1, and its locally ringed topos
extension, Remark 1.14, p.~380; the universal classifying topos has
quasi-compact final object. Apply this with the $S_n$ factor present,
so injectivity holds in the equivariant $K_0$ of the desired identity.

The pulled-back bundle has a filtration with line quotients
$L_1,\ldots,L_d$. Exterior-power and tensor-power filtrations permit
replacement by $\bigoplus_iL_i$. For this direct sum expand the
tensor product by its $n$ tensor positions. A multiplicity list
$(\nu_1,\ldots,\nu_d)$ with sum $n$ contributes
\[
 L_1^{\nu_1}\cdots L_d^{\nu_d}\otimes
 \operatorname{Ind}_{S_{\nu_1}\times\cdots\times S_{\nu_d}}^{S_n}1.
\]
The stabilizer acts trivially on tensor powers of a line. If
$\alpha_i$ counts indices with $\nu_j=i>0$, the induction class is
$\rho_1^{\alpha_1}\cdots\rho_n^{\alpha_n}$. Summing the distinct
multiplicity lists of this type gives the monomial symmetric function
$m_\mu(L_1,\ldots,L_d)$, where $\mu$ has $\alpha_i$ parts equal
to $i$. The fundamental theorem on symmetric polynomials writes this
as the integral polynomial $\Sigma_\alpha$ in the elementary classes
$\lambda^i(E)$. Injectivity of the flag pullback and then pullback
from the universal bundle prove the formula for $E$.

The right-hand side as a total series is also exponential under
addition. For finite alphabets of formal line elements its expansion is
\[
 \prod_j\left(\sum_{r\geq0}L_j^r\rho_r\right),\qquad\rho_0=1,
\]
so union of alphabets is multiplication. Degree by degree this is a
polynomial identity in elementary symmetric functions over $\mathbb Z$,
and hence holds for arbitrary pre-$\lambda$ inputs satisfying the
addition rule. Both sides extend identically to virtual classes. Thus
\[
 \tau_n(\xi)=\sum_{\sum i\alpha_i=n}
 \Sigma_\alpha(\lambda^1\xi,\ldots,\lambda^n\xi)
              \rho_1^{\alpha_1}\cdots\rho_n^{\alpha_n}.
\]
Bounded nonconstant rank is handled by its finite open-and-closed rank
decomposition. A-017 below justifies the infinite rank decomposition
as well, so the conclusion extends to all locally finite ranks.
\hfill$\square$

The universal polynomial $\Sigma_\alpha$ is homogeneous for weights
$\deg\lambda^i=i$. Its extreme terms are $m_{(1^n)}=e_n$ and
$m_{(n)}=p_n$; the coefficients of the formal monomials $\rho_1^n$
and $\rho_n$ are therefore $\lambda^n$ and the $n$th Newton operation.
If the target has relations among the $\rho_i$, these are coefficients
of the universal polynomial expression, not uniquely recoverable
coefficients of the target element. A-004 refutes the separate question
about effective direct summands.

\subsection*{A-011. Injection powers and the corrected fiber-type decomposition}
\textbf{Source status.} SA17, lines 498--515, describes cote 133 scans
20--23, which were not attached but have now been retrieved from the
official Montpellier facsimile and visually read (A-029). The
reported statements and the deletion on scan022 are verified against
those originals. For a finite $G$-set $E$, set
$\tau_n(E)=\operatorname{Inj}([n],E)$ and
$T_n(E)=\operatorname{Map}([n],E)$, with commuting $G$ and $S_n$ actions.
Then
\[
 \tau_n(E\amalg F)=\coprod_{i+j=n}
 \operatorname{Ind}_{G\times S_i\times S_j}^{G\times S_n}
       (\tau_i(E)\times\tau_j(F)).
\]
Indeed the inverse image of $E$ is an $i$-element subset of $[n]$;
choosing it is the coset space $S_n/(S_i\times S_j)$, and restrictions
of the injection give the remaining factors. This proves the
exponential property and its extension to $B(G)$ by group completion.

For a fiber type $\alpha=(\alpha_1,\ldots,\alpha_n)$ put
$m=\sum\alpha_i$, $n=\sum i\alpha_i$, and
$W_\alpha=\prod_i(S_i\wr S_{\alpha_i})$. The canonical surjection
$W_\alpha\to\prod_iS_{\alpha_i}\subset S_m$ permutes blocks of
each size. The summand of $T_n(E)$ of this type is
\[
 \operatorname{Ind}_{G\times W_\alpha}^{G\times S_n}
 \operatorname{Infl}_{G\times\prod_iS_{\alpha_i}}^{G\times W_\alpha}
 \operatorname{Res}_{G\times\prod_iS_{\alpha_i}}^{G\times S_m}
 \tau_m(E).
\]
\textbf{Proof.} A map determines its partition into nonempty fibers.
The stabilizer of a chosen partition of type $\alpha$ is $W_\alpha$.
Labeling its $m$ blocks by distinct elements of $E$ is an injection
from the block set, and the stabilizer acts through its permutation
quotient. Inducing over the orbit of the partition proves the formula.
The number of partitions of type $\alpha$ is
$n!/\prod_i\alpha_i!(i!)^{\alpha_i}$. Already for $n=2,m=1$,
induction from $S_1$ is wrong: diagonal maps have trivial $S_2$ action,
whereas induction from $S_1$ gives a free action. Thus the deletion
reported by SA17 is justified. \hfill$\square$

The operators $q_H^n$ were only named in the supplied study.
The original manuscript scan021 has now been retrieved and visually
read; A-029 gives its quotient definition, full extension proof and
precise polynomial-degree qualification. This resolves that earlier
source-availability gate.

\subsection*{A-012. Correcting the reported finite-injection-functor equivalence}
\textbf{Source status.} SA17, lines 517--532, describes cote133
scans24--27. These originally unattached pages have now been
visually read in the official facsimile (A-029); scan025 confirms
the single-map-functor assertion below, even for a sufficiently
large indexing set. Put
$h_I(E)=\operatorname{Inj}(I,E)$ and $H_J(E)=\operatorname{Map}(J,E)$
on finite sets and injections. The assertion that finite sums of
$h_I$ are equivalently direct summands of a single $H_J$ is false
as written.

\textbf{Proof and verification of the closure assertions.}
The functor $h_1\amalg h_1$ has cardinality two at a singleton,
whereas every $H_J$ has cardinality one there. It is neither a
coproduct summand nor a retract of a single $H_J$. The corrected
statement uses a finite sum of the $H_J$: each $h_I$ is a component
of $H_I$, consisting of its injective maps.

The category of elements of $h_I$ has initial object
$(I,\operatorname{id}_I)$, proving connectedness. The components
of $H_J$ are indexed by kernel equivalence relations on $J$; the
component with quotient $I$ is a copy of $h_I$. Products of
representables split into finitely many representables, according
to overlap equivalence relations of the two injection images.
By Yoneda, a natural map from $h_I$ to a finite sum of representables
is determined by a target component and an injection from its
representing set to $I$. The equalizer of two such maps is all
of $h_I$ if the choices agree, and empty otherwise. Thus finite
sums of representables are closed under finite limits, including
the terminal $h_\varnothing$.

An element of $h_I\circ h_J$ is an injective list indexed by $I$
of injections $J\to E$. Its equality pattern on $I\times J$ has
finitely many possibilities. Admissible patterns have no equalities
within a row and no two identical rows. Each pattern gives one
representable, by injecting its quotient set into $E$. Finite sums
in either argument are handled by the finitely many choices of
target component for each row. This proves closure under composition.

Finite sums of $h_I/H$ correspond to expressions
$\coprod_n h_{[n]}\times_{S_n}X_n$ with each $X_n$ finite and with
\emph{finite degree support}: decompose $X_n$ into $S_n$-orbits.
Dropping finite degree support would no longer describe a finite
sum. Eventual bijectivity satisfies two-out-of-three by taking the
maximum of the two thresholds; pullback preserves it pointwise.
Finally functors with injective transition maps are closed under
sums, products and composition, since each constituent transition
in the composite is an injection. These verify the substantive
closure assertions SA17 had grouped as ``definitions''.
\hfill$\square$

% The following block audits the separately reported scan 029 question.

\subsection*{A-013. Polynomial induction rings: fields, local rational algebras, lattices, and a singular counterexample}
\textbf{Source status.} SA17, lines 534--556, reports a question on
cote133 scan029, which was not attached but has now been visually
read in the official facsimile (A-029). Its local-ring question
is confirmed. In this paragraph $R_A(S_n)$
always means exact $K_0$ of representations finite projective over $A$.
For a local ring these underlying modules are finite free. The product
across $n$ is induction; $\rho_n$ is the trivial representation in degree
$n$. This is different from the split group used in A-005.

Over a characteristic-zero field the induction ring is the symmetric
function ring $\mathbb Z[h_1,h_2,\ldots]$, with $\rho_i$ corresponding to
$h_i$. This is the classical Frobenius characteristic theorem (Macdonald, second edition, I, (7.3), p.~113, and I, (2.8), p.~22). The same
conclusion holds for every local $\mathbb Q$-algebra $A$: the split
algebra $\mathbb Q[S_n]=\prod_{\lambda\vdash n}\operatorname{Mat}_{f^\lambda}(\mathbb Q)$
base-changes to the analogous matrix-algebra product over $A$.
Morita equivalence identifies the representation exact category with a
product of finite projective $A$-module categories. Such modules are
free because $A$ is local. The rational Specht bases, and the induction
matrices between them, therefore base-change unchanged.

For $A=\mathbb F_2$, the regular $S_2$ module has a trivial submodule
and trivial quotient, giving $\rho_1^2=2\rho_2$. Thus the polynomial
assertion fails. More generally $R_{\mathbb F_p}(S_p)$ has rank
$p(p)-1$, since simples are indexed by $p$-regular partitions and
$(1^p)$ is the unique excluded partition of $p$ (James, \emph{The
Representation Theory of the Symmetric Groups}, Theorem 11.5, p.~40 (with Definition 11.2, p.~39)). A polynomial ring with one generator in each
positive degree would have rank $p(p)$ in this degree.

Characteristic zero alone is also insufficient for local rings.
Let
\[
 A=\mathbb Z_{(2)}[\epsilon]/(2\epsilon,\epsilon^2).
\]
This is a local ring of characteristic zero, with nonzero $\epsilon$.
The unit $1+\epsilon$ has square one, so defines a rank-one free
$A[S_2]$ module $L$ on which the transposition acts by $1+\epsilon$.
Trace of this transposition is additive in exact sequences of
$A$-free representations, since these sequences split as $A$-modules.
It therefore defines a homomorphism $R_A(S_2)\to A$. Its values on
$\rho_1^2$ and $\rho_2$ are $0$ and $1$, while its value on $L$ is
$1+\epsilon$. The latter is not in the integer subgroup of $A$.
Thus $[L]$ is not in the degree-two image of
$\mathbb Z[\rho_1,\rho_2,\ldots]$, proving failure of surjectivity.

\paragraph{The specific $\mathbb Z_p$-lattice question has an affirmative answer.}
Let $O=\mathbb Z_p$, $K=\mathbb Q_p$, and $R=O[S_n]$. The generic-fiber
map
\[
 K_0(\text{$O$-free finite $R$-modules})\longrightarrow R_K(S_n)
\]
is an isomorphism. Here is a proof spelling out the torsion issue.
Let $G_0(R)$ use all finitely generated $R$-modules. Any such module
$M$ has a finite free $R$-module cover $P\to M$ whose kernel $L$ is
$O$-torsion-free, hence $O$-free of finite rank. Set
$\chi(M)=[P]-[L]$ in the lattice $K_0$. Schanuel's lemma shows
independence of the cover. For a short exact sequence, choose free
covers of the outside terms and lift the quotient cover through the
middle term. The resulting direct-sum cover has kernels in a short
exact sequence; this proves additivity of $\chi$. On a lattice it
is its original class. Hence the lattice $K_0$ equals $G_0(R)$.

Every $O$-torsion $R$-module has a filtration by powers of $p$, with
quotients that are finite $\mathbb F_p[S_n]$ modules. The classes of
these modular modules are integer combinations of reductions of
Specht lattices: the submatrix indexed by $p$-regular partitions is
unitriangular (James, Corollary 12.3, p.~43; compatibility with
integral reduction is Corollary 8.11, p.~31). For each Specht lattice
$S_O^\lambda$, the exact sequence
\[
 0\longrightarrow S_O^\lambda\xrightarrow{p}S_O^\lambda
 \longrightarrow S_O^\lambda/pS_O^\lambda\longrightarrow0
\]
shows that its reduction is zero in $G_0(R)$. Therefore all torsion
module classes are zero there.

Any $K[S_n]$ module has an $S_n$-stable $O$-lattice, obtained by
summing translates of an arbitrary full lattice, proving surjectivity.
If two lattice classes have the same generic class, semisimplicity
over $K$ makes their generic representations isomorphic. Multiplying
such an isomorphism by a power of $p$ embeds one lattice into the
other with torsion cokernel. That cokernel is zero in $G_0(R)$, so
the lattice classes agree. This proves injectivity. Generic fiber
commutes with induction, hence
\[
 \bigoplus_n R_{\mathbb Z_p}(S_n)\simeq
 \bigoplus_n R_{\mathbb Q_p}(S_n)=\mathbb Z[\rho_1,\rho_2,\ldots].
\]
The argument is specific to these exact lattice categories and the
surjective decomposition map for symmetric groups. It does not extend
to arbitrary singular local rings, as the explicit $A$ above shows.
A classification of all local coefficient rings is not asserted.


\subsection*{A-014. The big-Witt calculations of cote 43}
\textbf{Source.} Cote 43, scan 037, lines 3--50. These are statements
about the big Witt ring in exterior-series convention. They do not
describe subset operations in a Burnside ring. With this convention
the calculations are valid, with continuity and naturality understood
in the final uniqueness assertion.

\textbf{Verification.} Put $W^+(A)=1+tA[[t]]$, with addition given
by ordinary multiplication of series. On formal roots its other
product is
\[
 \prod_i(1+x_it)*\prod_j(1+y_jt)=\prod_{i,j}(1+x_iy_jt).
\]
Each coefficient is symmetric in both alphabets and therefore is
an integral polynomial in their elementary symmetric functions.
These polynomials define the product for every ring and give an
affine ring scheme over $\mathbb Z$ with underlying scheme
$\operatorname{Spec}\mathbb Z[a_1,a_2,\ldots]$. Associativity and
distributivity are identities on formal roots, hence polynomial
identities over $\mathbb Z$. The standard special $\lambda$
structure is that of SGA 6, V, 2.3--2.5, especially Example 2.5(b), pp.~321--322; the line identity is
\[
 (1+at)*f(t)=f(at).
\]
If $A$ has characteristic $p$, coefficientwise Frobenius
$F(\sum a_it^i)=\sum a_i^pt^i$ is a ring endomorphism because
the operations are integral polynomials. The map $V(f)(t)=f(t^p)$
is injective and additive for Witt addition. The ordinary series
identity $f(t)^p=F(f)(t^p)$ proves
\[
 FV=VF=p
\]
as endomorphisms of the additive Witt group. If $x=a^{p^i}$, then
$1+xt^{p^i}=(1+at)^{p^i}$ and distributivity plus the line identity
give
\[
 (1+xt^{p^i})*f(t)=F^i(f)(xt^{p^i}).
\]
For arbitrary $x$, adjoin a root of $T^{p^i}-x$. This is a free,
faithfully flat extension, so the identity descends. This verifies
the characteristic-$p$ formulas without a perfectness assumption.

For the last map to $p$-typical Witt vectors, convert the exterior
series to the usual convention by $f(t)\mapsto f(-t)^{-1}$ and
take the usual $p$-typical projection, retaining the ghost coordinates
indexed by $1,p,p^2,\ldots$. Equivalently, on a finite formal root
alphabet set
\[
 \psi\left(\prod_i(1+x_it)\right)=\sum_i[x_i]\in W_p(A).
\]
The coordinates of the right side are symmetric integral
polynomials in the $x_i$, using the defining integral addition
polynomials of the Witt functor; its $j$th coordinate has degree
$p^j$, so depends on only finitely many coefficients of the input
series. This constructs a continuous natural map on all series.
The ghost identities show that it respects both ring operations
over the universal torsion-free polynomial ring, and hence, as
integral polynomial identities, over every ring. It sends
$1+at$ to the Teichmuller vector $(a,0,\ldots)$. A continuous
natural homomorphism with these line values is unique: split
universal finite polynomials into formal roots, use additivity,
and pass coefficientwise to formal series. Without this natural
continuous interpretation, the assertion is not a claim that the
line elements generate $W^+(A)$ as an abstract abelian group.
\hfill$\square$

\subsection*{A-015. Tensor functor identities and the finite-operator viewpoint}
\textbf{Source.} Cote 2, scan 017, lines 2--38; cote 12, scan 041,
lines 22--24; cote 119, scan 217, lines 25--31.
For group representations on finite locally free modules,
restriction preserves tensor products and exterior powers. External
tensor product is the tensor of the two pullbacks; pullback by the
diagonal gives the internal tensor product. These identities follow
directly from the actions on elementary tensors. Restrictions and
inductions are transitive, by the associativity isomorphism for
tensoring bimodules.

The projection formula for induction along $u:G\to H$ is the explicit
isomorphism
\[
 kH\otimes_{kG}(M\otimes\operatorname{Res}N)
 \longrightarrow(kH\otimes_{kG}M)\otimes N,
 \qquad h\otimes(m\otimes n)\longmapsto(h\otimes m)\otimes hn.
\]
Its inverse sends $(h\otimes m)\otimes n$ to
$h\otimes(m\otimes h^{-1}n)$; both maps respect the balancing
relations and the $H$ action. For finite groups these constructions
preserve finite locally free underlying modules in the usual
subgroup-induction setting. If $u_*$ means right-adjoint
coinduction instead, the projection formula holds when $N$ is
dualizable, by tensor-Hom adjunction. It is not asserted here for
arbitrary sheaves without this hypothesis.

The statement on cote 12 is a programme to determine $\lambda$
operations using finite operator groups, not a positivity theorem.
The historical remark in cote 119 supplies no further mathematical
identity. The precise familiar dictionary is that, if $r!$ is
invertible, $\bigwedge^r E$ is the image of the antisymmetrizer
$r!^{-1}\sum_{\pi\in S_r}\operatorname{sgn}(\pi)\pi$ on
$E^{\otimes r}$. Over $\mathbb Z$ the exterior quotient and the
exterior filtrations still exist, but this averaging idempotent
need not. SGA 6, V, Example 2.2(b) and Lemma 2.2.1, pp.~315--316,
give the exact-category construction. None of these formal
identities alone provides positive irreducible multiplicities.
\hfill$\square$

For the classifying-topos reduction on cote 133 scan 006, lines
34--45, add the local-standard-form hypothesis explicitly. If
$q_m,q_n$ are the chosen standard forms and the given forms are
locally isometric to them, their sheaves of isometries are
$O(q_m)$ and $O(q_n)$ torsors. The associated-bundle functor from
representations of their product preserves orthogonal sums,
tensor products, and exterior powers. It therefore pulls any
universal identity for the standard forms back to the given pair;
one factor suffices for an iteration identity. In the usual etale
setting with two invertible, nonsingular symmetric forms are
locally standard. An arbitrary ringed topos need not already have
that topology, so this step is conditional on the stated local
triviality, or must be formulated on the universal stack of
quadratic forms. A-005 still refutes specialness for the split
group in the archive's generality.

\subsection*{A-016. What the supplied Borger source does and does not transfer}
In the supplied \texttt{lrfoe13.tex}, \S1.1 (lines 774--807)
describes $\Lambda$-structures and, for flat schemes over
$\mathbb Z$, commuting Frobenius lifts. These are special
$\lambda$ structures. A-002 therefore blocks an automatic transfer
from the canonical Burnside pre-structure to that geometry.
\S2.5 (lines 956--977) gives the Chebychev example
$\psi_2(x)=x^2-2$ on the representation ring of $SL_2$;
the existence of a $\Lambda$-structure does not assert preservation
of an arbitrary coefficientwise positive cone. Its cohomological
theorem in \S6, labelled
\texttt{thm:lambda-varieties-are-abelian} (lines 1863--1910), and
the following corollary
\texttt{cor:lambda-varieties-are-Artin-Tate} (1913--1944), require
an actual finite-type $\Lambda$-scheme. They supply no such scheme
with fibers or cohomology equal to the required Kronecker
multiplicity spaces. That missing realization, including its
grading and a proof of support saturation, cannot be replaced
by identifying the coefficient ring as a $\lambda$-ring.

The section titled ``Relation to finite sets'', starting at line
2682, occurs after the first \verb|\end{document}| at line 2680.
It is inactive source material in this supplied draft; it must
not be cited as a numbered theorem of the compiled paper. The
source also describes itself as preliminary at lines 441--443.
These observations restrict the claimed dictionary; they do not
refute Borger's stated theorems or assert a general impossibility
of a future geometric Kronecker rule.

\subsection*{A-017. Removing the global rank bound in A-010}
\textbf{Claim.} If a locally ringed topos decomposes into open-and-closed
rank strata, the exact $K_0$ of its category of locally free modules
of finite local rank is the product of the corresponding exact
$K_0$ groups of the strata. Consequently A-010 holds for virtual
locally free classes without a global rank bound.

\textbf{Proof.} We prove the elementary exact-category statement
\[
 K_0\left(\prod_i\mathcal C_i\right)\simeq\prod_iK_0(\mathcal C_i).
\]
Every element of an exact $K_0$ is a difference of two object classes;
choosing these coordinatewise proves surjectivity. For injectivity
write an element as $[A]-[B]$ and suppose every coordinate is zero.
In the split group completion of $\mathcal C_i$, the difference
$[A_i]-[B_i]$ is a finite signed sum of short-exact-sequence
relations. Combine all positive relations by direct sum into
one sequence $U_i\to V_i\to W_i$, and all negative relations into
one sequence $U_i'\to V_i'\to W_i'$. Thus
\[
 [A_i]-[B_i]=[V_i]-[U_i]-[W_i]-[V_i']+[U_i']+[W_i']
\]
in the split group completion. By the defining equivalence of a
commutative-monoid group completion there is an object $D_i$ and
an isomorphism
\[
 A_i\oplus U_i\oplus W_i\oplus V_i'\oplus D_i
 \simeq B_i\oplus V_i\oplus U_i'\oplus W_i'\oplus D_i.
\]
The coordinate tuples form objects and exact sequences in the
product category. Their isomorphism, together with its two exact
relations, proves $[A]=[B]$ there. No uniform bound on the number
of original relations is needed after folding them into two
sequences. This proves the product assertion.

Now write a virtual class as $[E]-[F]$ and decompose the topos by
the pair of locally constant ranks of $E,F$. Modules on this
disjoint decomposition glue coordinatewise, so the exact category
of locally free modules of finite local rank is the product of
the categories on the strata. The same holds after adjoining the
finite $S_n$ action. The two sides of A-010 agree on each stratum
by its constant-rank proof, hence agree globally by the product
assertion. This establishes the claimed extension. \hfill$\square$

\subsection*{A-018. A positive kernel model for every fixed-height Kronecker family}
\textbf{Statement.} Let $\lambda,\mu,\nu\vdash N$ have at most
$p,q,r$ parts, respectively, and pad them by zeros to these lengths.
Let $\mathcal T(\lambda,\mu,\nu)$ be the set of nonnegative integer
$p\times q\times r$ arrays $a=(a_{ijk})$ with one-dimensional margins
$\lambda,\mu,\nu$. For $1\leq i<p$ define the integer matrix
\[
 D_i^{(1)} e_a=\sum_{j=1}^q\sum_{k=1}^r a_{i+1,j,k}
 e_{a-e_{i+1,j,k}+e_{i,j,k}},
\]
from the vector space with basis $\mathcal T(\lambda,\mu,\nu)$ to
that with margins $(\lambda+\epsilon_i-\epsilon_{i+1},\mu,\nu)$.
Terms with zero coefficient are omitted. Define $D_j^{(2)}$ and
$D_k^{(3)}$ by the identical rule on the other axes. The target
margins need not be partitions; an impossible margin has zero
target space. If $D$ is the vertical stack of all these matrices,
then
\[
 \boxed{\quad g(\lambda,\mu,\nu)=\dim_{\mathbb Q}\ker D
 =|\mathcal T(\lambda,\mu,\nu)|-\operatorname{rank}_{\mathbb Q}D.\quad}
\]
This gives explicit multiplicity spaces and an exact algorithm
polynomial in $N$ for fixed $p,q,r$. It does not assert polynomial
complexity in the binary encoding length of $N$, or when heights
vary. The binary-height instance is F-001; this statement includes
all fixed-height triples and fixed-height rectangular triples.

\textbf{Proof.} Take $U=\mathbb Q^p$, $V=\mathbb Q^q$,
$W=\mathbb Q^r$. In $\operatorname{Sym}^N(U\otimes V\otimes W)$
the monomials $\prod x_{ijk}^{a_{ijk}}$ form a basis, and the three
diagonal tori have weight equal to the three margins. Thus the
weight $(\lambda,\mu,\nu)$ space has exactly the indicated basis.
The elementary upper-triangular Lie algebra operator
$E_{i,i+1}$ on $U$ acts as
$\sum_{j,k}x_{ijk}\partial/\partial x_{i+1,j,k}$, giving precisely
$D_i^{(1)}$. The same computation gives the other two sets of
matrices. Simple positive root operators generate all positive
root operators by commutators. In characteristic zero their
common kernel is the subspace fixed by the upper unipotent group:
on a finite-dimensional polynomial representation each root
subgroup acts as the finite exponential of its nilpotent
derivative.

For clarity the multiplicity calculation is a finite-group
calculation. Reordering tensor factors and taking diagonal
$S_N$ invariants gives
\[
 \operatorname{Sym}^N(U\otimes V\otimes W)
 \simeq (U^{\otimes N}\otimes V^{\otimes N}
                  \otimes W^{\otimes N})^{S_N}.
\]
Characteristic-zero averaging identifies symmetric coinvariants
with these invariants. Apply Schur--Weyl decomposition separately
to the three tensor factors. This yields
\[
 \bigoplus_{\alpha,\beta,\gamma\vdash N}
 S_\alpha U\otimes S_\beta V\otimes S_\gamma W\otimes
 ([\alpha]\otimes[\beta]\otimes[\gamma])^{S_N}.
\]
Terms whose partition has too many rows vanish. Symmetric-group
irreducibles are self-dual, so the last space has dimension
$g(\alpha,\beta,\gamma)$. In the irreducible polynomial
$GL(U)$ module $S_\alpha U$, the simultaneous positive-root
kernel is its one-dimensional highest-weight line of weight
$\alpha$. One may check uniqueness directly: order products of
matrix Lie algebra generators using their commutator relations,
with lowering operators first, diagonal operators next, and
raising operators last. A weight vector killed by all raising
operators generates only weights below its own. Irreducibility
then forces its weight to be the highest weight $\alpha$, whose
weight space has dimension one. The same applies to $V,W$.
Consequently the joint kernel in weight $(\lambda,\mu,\nu)$ is
exactly the corresponding invariant multiplicity space. The
integer matrices have the same rank over $\mathbb Q$ and
$\mathbb C$, proving the formula. Schur--Weyl here is the standard
characteristic-zero equivalence in Macdonald, Chapter I,
Appendix A, (5.3)--(5.6), pp.~155--156.

There are
$D_0=pqr-p-q-r+2$ free table entries: after choosing all cells
except $(1,1,1)$, $(i,1,1)$ for $i>1$, $(1,j,1)$ for $j>1$,
and $(1,1,k)$ for $k>1$, the margins determine those remaining
cells uniquely. Each chosen entry lies between zero and $N$;
nonnegative completion can be checked directly. Thus every
source and target table set has at most $(N+1)^{D_0}$ elements.
For fixed $p,q,r$ the number of targets and the number of nonzero
entries per column are bounded independently of $N$, and matrix
entries have magnitude at most $N$. Fraction-free integer
elimination has polynomial bit complexity in these dimensions
and in $\log(N+1)$: intermediate determinants have bit lengths
bounded by the Hadamard bound, polynomial in the matrix dimension
and entry bit length. This proves the stated polynomial-in-$N$
algorithm, including zero and degenerate margins. \hfill$\square$

\textbf{Geometric and positivity scope.} Equivalently this is the
weight $(\lambda,\mu,\nu)$ part of
$H^0(\mathfrak u_p\oplus\mathfrak u_q\oplus\mathfrak u_r,
\operatorname{Sym}^N(U\otimes V\otimes W))$.
All spaces and matrices are explicitly defined over $\mathbb Z$,
using the coordinate affine scheme and its weight grading; the
diagonalizable-group/graded-module dictionary is SGA 3, I,
Proposition 4.7.3, pp.~35--36. The equality with ordinary
Kronecker coefficients is asserted only over characteristic zero.
This is a constructive positive vector-space realization and
finite rank algorithm, not a cancellation-free integer-point
count, a saturation theorem, or a uniformly efficient rule when
the determinant size and rectangle heights grow.

\paragraph{A direct finite-$G$-set construction of the same kernel.}
This model can be placed before linearization. Let $G=S_N$ and
let $X_\lambda$ be the finite $G$-set of ordered decompositions
$[N]=B_1\sqcup\cdots\sqcup B_p$ with $|B_i|=\lambda_i$.
For each $i<p$, the correspondence which moves one element
from $B_{i+1}$ to $B_i$ gives the equivariant linear map
\[
 \partial_i:\mathbb Q[X_\lambda]\longrightarrow
                 \mathbb Q[X_{\lambda+\epsilon_i-\epsilon_{i+1}}],
 \qquad [B]\longmapsto
                 \sum_{u\in B_{i+1}}[B\text{ with }u\text{ moved}].
\]
Indeed, identify $\mathbb Q[X_\lambda]$ with the weight-$\lambda$
subspace of $(\mathbb Q^p)^{\otimes N}$ by recording the block
of every tensor position. Then $\partial_i$ is exactly the
simple positive-root operator. Schur--Weyl and the highest-weight
argument above identify its joint kernel with the Specht module
$[\lambda]$, including its $S_N$ action.

Take the three corresponding systems for $\lambda,\mu,\nu$.
Tensoring over a field preserves kernels, and invariants preserve
kernels; hence the joint kernel inside
\[
 \bigl(\mathbb Q[X_\lambda]\otimes\mathbb Q[X_\mu]
                     \otimes\mathbb Q[X_\nu]\bigr)^{S_N}
\]
is $([\lambda]\otimes[\mu]\otimes[\nu])^{S_N}$.
The orbits of triples of ordered decompositions are indexed by
the intersection tables
$a_{ijk}=|B_i\cap C_j\cap E_k|$, giving precisely
$\mathcal T(\lambda,\mu,\nu)$. Use the normalized orbit sums
as a rational invariant basis. From any point of orbit $a$,
there are $a_{i+1,j,k}$ ways to move an element in the indicated
intersection to the new block. Equivariance distributes those
edges uniformly over the target orbit. Therefore the coefficient
on its normalized orbit sum is exactly $a_{i+1,j,k}$, which
recovers the matrix $D_i^{(1)}$, and likewise on the other axes.
This proves the formula directly through finite $G$-sets and
their incidence correspondences. It uses actual kernels after
linearization, without claiming that a nontrivial irreducible
itself has an effective permutation lift. The constant-group
scheme dictionary is SGA 3, I, 4.1, p.~23; the representation
comodule dictionary is I, Proposition 4.7.2, pp.~34--35.

\subsection*{A-019. A sufficient positive-basis hypothesis and its exact saturation scope}
\textbf{Statement.} Let $A=\bigoplus_{\alpha\in L_0}A_\alpha$ be
a commutative domain graded by a finite-rank free abelian group $L_0$, over
an algebraically closed field $k$. Suppose $A$ is integrally
closed in its fraction field and each $A_\alpha$ has dimension
at most one over $k$. Put
$S=\{\alpha:A_\alpha\ne0\}$ and $L=\mathbb ZS\subseteq L_0$.
Then $S$ is saturated in $L$: if $\alpha\in L$ and
$h\alpha\in S$ for an integer $h\geq1$, then $\alpha\in S$.
Moreover there is a homogeneous basis $(u_\alpha)_{\alpha\in S}$
with $u_\alpha u_\beta=u_{\alpha+\beta}$; thus
$A\simeq k[S]$. Finite generation is not needed for these
assertions. If $A$ is finitely generated, this supplies an affine
semigroup algebra with an effective monomial basis. If in
addition $L$ is primitive in $L_0$, the same saturation implication
holds for every $\alpha\in L_0$.

\textbf{Proof of saturation.} The support is a monoid because
$A$ is a domain. Write $\alpha=\beta-\gamma$ with
$\beta,\gamma\in S$; every element of $\mathbb ZS$ has such an
expression by collecting the positive and negative summands.
Choose nonzero $f_\beta\in A_\beta$,
$f_\gamma\in A_\gamma$, and $f_{h\alpha}\in A_{h\alpha}$.
Both $f_\beta^h$ and $f_\gamma^h f_{h\alpha}$ are nonzero
elements of the same at-most-one-dimensional space $A_{h\beta}$.
Hence $z=f_\beta/f_\gamma$ satisfies
$z^h=c f_{h\alpha}$ for some $c\in k^\times$.
It is integral over $A$ and therefore lies in $A$ by normality.
Its degree is $\alpha$: in the homogeneous decomposition of
$z$, multiplying by the nonzero homogeneous $f_\gamma$ can
produce the single degree $\beta$ only from degree $\alpha$.
Thus $A_\alpha\ne0$. If $L$ is primitive and
$h\alpha\in S\subseteq L$ for $\alpha\in L_0$, then
$\alpha\in L$ by primitivity, proving the last assertion.

\textbf{Proof of the positive basis.} Localize $A$ by all its
nonzero homogeneous elements, obtaining a graded domain $B$.
Its degree-zero part is $k$: the numerator and denominator of a
homogeneous fraction of degree zero belong to a common
one-dimensional $A$-piece and hence are proportional. Every
degree in $L$ occurs, and every nonzero homogeneous fraction is
invertible. Consequently each $B_\alpha$ is one-dimensional.
Choose a $\mathbb Z$-basis $\ell_1,\ldots,\ell_s$ of $L$ and
a nonzero $z_i\in B_{\ell_i}$. Distinct Laurent monomials in
the $z_i$ have distinct degrees, so they are linearly independent.
They span $B$ because every homogeneous piece is a line. Hence
$B=k[z_1^{\pm1},\ldots,z_s^{\pm1}]$ as a graded algebra.
For $\alpha\in S$, the corresponding Laurent monomial
$u_\alpha$ spans $B_\alpha$ and belongs to $A_\alpha$, since
that is a nonzero line in $B_\alpha$. These monomials give the
claimed multiplicative basis and algebra isomorphism.
\hfill$\square$

The algebraic-closure assumption is harmless for the intended
complex representation application; the displayed proof in
fact needs only that $A_0=k$ and the stated one-dimensional
grading. It is explicit multiplication plus normality and a
multiplicity-free grading that produce the positive saturated
model. An abstract special $\lambda$-ring does not provide these
hypotheses. In a highest-weight algebra whose graded dimensions
are Kronecker coefficients, the one-dimensional hypothesis fails
as soon as some coefficient exceeds one; no global application
to that algebra is asserted here.

The lattice qualification cannot be dropped. The polynomial
algebra $k[t]$ with $\deg t=2$ is normal and multiplicity-free,
but its support $2\mathbb N$ is saturated only in $2\mathbb Z$,
not in the chosen ambient lattice $\mathbb Z$. Nor does a
saturated finer support guarantee a saturated projected grading:
$\mathbb N^2$ is saturated in $\mathbb Z^2$, while its image under
$(a,b)\mapsto2a+3b$ is $\langle2,3\rangle$, which misses $1$
although it generates all of $\mathbb Z$ and contains $2$.
A proposed toric degeneration must therefore also prove
saturation of the particular weight projection that records
the desired multiplicities.
